\section*{Capítulo \ref{ch:b1:eukaryotic-cell} --- Organización funcional de la célula eucariota}

\begin{solution}{exo:b1:eukaryotic-cell:1}
\omterm{def:b1:eukaryotic-cell:nucleus}{Núcleo} (cromosomas, transcripción); RE rugoso (proteínas de secreción y de
membrana); RE liso (lípidos, calcio, destoxificación); \omterm{def:b1:eukaryotic-cell:endomembrane}{Golgi} (clasificación,
modificación, empaquetado); \omterm{def:b1:eukaryotic-cell:endomembrane}{lisosomas} (digestión); \omterm{def:b1:eukaryotic-cell:peroxisome}{peroxisomas} (oxidaciones,
catalasa); \omterm{def:b1:eukaryotic-cell:mitochondrion}{mitocondrias} (respiración, ATP).
\end{solution}

\begin{solution}{exo:b1:eukaryotic-cell:2}
La sintetizan ribosomas unidos al RE rugoso y se enhebra en su luz
(\omterm{prop:b1:eukaryotic-cell:secretory}{secuencia señal}); del RE brotan vesículas que se fusionan con el \omterm{def:b1:eukaryotic-cell:endomembrane}{Golgi} cis;
la proteína cruza la pila, se modifica y se clasifica; abandona la cara
trans en una \omterm{def:b1:eukaryotic-cell:plantcell}{vacuola} de condensación que madura hasta gránulo de secreción;
el gránulo se fusiona con la membrana apical ante una señal y se vacía en el
conducto.
\end{solution}

\begin{solution}{exo:b1:eukaryotic-cell:3}
$7 + 32 = \qty{39}{\%}$. La membrana interna está plegada en \omterm{def:b1:eukaryotic-cell:mitochondrion}{crestas} porque
alberga la cadena respiratoria y la ATP sintasa, cuyo rendimiento es
proporcional a su área; la externa es una simple envoltura.
\end{solution}

\begin{solution}{exo:b1:eukaryotic-cell:4}
Vegetal: pared de celulosa, \omterm{def:b1:eukaryotic-cell:plantcell}{vacuola} grande, \omterm{def:b1:eukaryotic-cell:plastid}{plastos} (y también
\omterm{def:b1:eukaryotic-cell:plantcell}{plasmodesmos}). Animal: \omterm{def:b1:eukaryotic-cell:endomembrane}{lisosomas}, centriolos (y la capacidad de reptar y
fagocitar).
\end{solution}

\begin{solution}{exo:b1:eukaryotic-cell:5}
El RE a los \qty{3}{min} (ya en descenso), el \omterm{def:b1:eukaryotic-cell:endomembrane}{Golgi} hacia los \qty{15}{min}
y los gránulos hacia los \qty{60}{min}. Tránsito del RE al gránulo, unos
\qty{45}{min}; la mitad secretada hacia los \qty{110}{min}.
\end{solution}

\begin{solution}{exo:b1:eukaryotic-cell:6}
El \omterm{def:b1:eukaryotic-cell:endomembrane}{Golgi}, mantenido cerca del \omterm{def:b1:eukaryotic-cell:cytoskeleton}{centrosoma} por la dineína sobre los
\omterm{def:b1:eukaryotic-cell:cytoskeleton}{microtúbulos}, se fragmenta y se dispersa; el tráfico de vesículas se reduce
a la difusión y el transporte a la periferia se detiene; el huso no puede
formarse y las \omterm{def:b1:cell-unit-of-life:cell}{células} se detienen en mitosis; los cilios, cuyo eje son
\omterm{def:b1:eukaryotic-cell:cytoskeleton}{microtúbulos}, se paran (los ya existentes son estables, pero no se
renuevan).
\end{solution}

\begin{solution}{exo:b1:eukaryotic-cell:7}
Dos membranas, la interna de tipo bacteriano; ADN circular sin histonas;
ribosomas de tamaño bacteriano y sensibles a los antibióticos; división por
fisión; secuencias génicas que se agrupan con las alfaproteobacterias. No
puede vivir sola porque la mayoría de sus genes se han trasladado al \omterm{def:b1:eukaryotic-cell:nucleus}{núcleo}:
importa del \omterm{def:b1:eukaryotic-cell:organelle}{citosol} casi todas sus proteínas.
\end{solution}

\begin{solution}{exo:b1:eukaryotic-cell:8}
Los azúcares se añaden en la cara luminal en el RE y en el \omterm{def:b1:eukaryotic-cell:endomembrane}{Golgi}. Cada
fusión de vesículas conserva la orientación: la cara luminal de una vesícula
pasa a ser la cara luminal del compartimento siguiente, y cuando la vesícula
se fusiona con la membrana plasmática su luz se abre al exterior, de modo
que los azúcares miran hacia fuera. Lo que es luminal está, topológicamente,
fuera.
\end{solution}

\begin{solution}{exo:b1:eukaryotic-cell:9}
Cada gránulo, $\pi d^2 = \qty{3.14}{\micro m^2}$; al día,
\qty{6280}{\micro m^2}, es decir, \qty{157}{} veces la cara apical. La
\omterm{def:b1:cell-unit-of-life:cell}{célula} debe recuperar membrana por \omterm{def:b1:eukaryotic-cell:endocytosis}{endocitosis} al mismo ritmo, de modo que
la membrana apical se renueva cada nueve minutos y la membrana recuperada
vuelve al \omterm{def:b1:eukaryotic-cell:endomembrane}{Golgi} para reutilizarse.
\end{solution}

\begin{solution}{exo:b1:eukaryotic-cell:10}
La cadena de \qty{320}{} residuos es el precursor con su \omterm{prop:b1:eukaryotic-cell:secretory}{secuencia señal} de
\qty{20}{} residuos, que normalmente se corta al entrar en el RE. Con
vesículas, la cadena entra en la luz durante la síntesis, la señal se corta
y la membrana protege al producto de la proteasa. El mutante sin la señal no
lo reconoce la maquinaria del RE: se fabrica en ribosomas libres y se queda
en el \omterm{def:b1:eukaryotic-cell:organelle}{citosol}, con la longitud completa, sin secretarse nunca.
\end{solution}

\begin{solution}{exo:b1:eukaryotic-cell:11}
Las membranas internalizadas por \omterm{def:b1:eukaryotic-cell:endocytosis}{endocitosis} y los \omterm{def:b1:cell-unit-of-life:prokeuk}{orgánulos} gastados se
llevan a los \omterm{def:b1:eukaryotic-cell:endomembrane}{lisosomas} para digerirlos; sin la enzima, el lípido que
contienen no puede degradarse y llena los \omterm{def:b1:eukaryotic-cell:endomembrane}{lisosomas}, que se hinchan hasta
que la \omterm{def:b1:cell-unit-of-life:cell}{célula} falla. Las \omterm{def:b1:cell-unit-of-life:cell}{células} con el recambio de membrana más rápido (las
\omterm{def:b1:body-plans-tissues:nervous}{neuronas}, que reciclan sin cesar vesículas sinápticas; los macrófagos) lo
acumulan antes. Una enzima inyectada que lleve la señal de azúcar propia de
las enzimas lisosómicas puede unirse a receptores de la superficie celular y
ser captada por \omterm{def:b1:eukaryotic-cell:endocytosis}{endocitosis} mediada por receptor hasta los \omterm{def:b1:eukaryotic-cell:endomembrane}{lisosomas}: es el
principio de la terapia de sustitución enzimática.
\end{solution}

\begin{solution}{exo:b1:eukaryotic-cell:12}
La endosimbiosis da buena cuenta de las \omterm{def:b1:eukaryotic-cell:mitochondrion}{mitocondrias} y los \omterm{def:b1:eukaryotic-cell:plastid}{plastos} (dos
membranas, genoma y ribosomas propios, fisión, parientes bacterianos) y, por
tanto, del metabolismo energético del eucariota. No explica el \omterm{def:b1:eukaryotic-cell:nucleus}{núcleo}, el RE
y el \omterm{def:b1:eukaryotic-cell:endomembrane}{Golgi}, ni el \omterm{def:b1:eukaryotic-cell:cytoskeleton}{citoesqueleto}, que no tienen equivalentes bacterianos y
parecen invenciones del linaje hospedador: repliegues de su membrana y
proteínas nuevas. ``Federación'' es acertado para los \omterm{def:b1:cell-unit-of-life:prokeuk}{orgánulos}; ``mantenida
unida por un sistema de membranas'' apunta a la parte mayor y sin explicar:
una \omterm{def:b1:cell-unit-of-life:cell}{célula} hospedadora que ya era eucariota en su arquitectura antes de
acoger a las bacterias.
\end{solution}

\begin{solution}{pb:b1:eukaryotic-cell:1}
\textbf{1.} Un pulso corto marca una cohorte estrecha de moléculas, de modo
que su posición señala un instante; el exceso de aminoácido frío detiene la
incorporación posterior, así que la cohorte no se difumina por síntesis
posterior.
\textbf{2.} RE: \qty{3}{min} (el punto más temprano); \omterm{def:b1:eukaryotic-cell:endomembrane}{Golgi} y \omterm{def:b1:eukaryotic-cell:plantcell}{vacuolas} de
condensación: hacia los \qty{15}{min}; gránulos: hacia los \qty{60}{min}.
\textbf{3.} Del RE al \omterm{def:b1:eukaryotic-cell:endomembrane}{Golgi}, unos \qty{10}{min}; del \omterm{def:b1:eukaryotic-cell:endomembrane}{Golgi} al gránulo, unos
\qty{45}{min}.
\textbf{4.} Sobre todo en gránulos (\qty{70}{\%}), una décima parte en el
\omterm{def:b1:eukaryotic-cell:endomembrane}{Golgi} y otra décima ya secretada: la vía constitutiva y los primeros
gránulos formados liberan algo, mientras que la mayor parte espera un
estímulo.
\textbf{5.} La gemación y el transporte de vesículas exigen energía y
membranas fluidas; con frío se detiene la etapa del RE al \omterm{def:b1:eukaryotic-cell:endomembrane}{Golgi} y la
proteína se acumula aguas arriba del bloqueo.
\textbf{6.} $10\,\mathrm{g}/\num{8e10} = \qty{1.25e-10}{g} =
\qty{125}{pg}$ por \omterm{def:b1:cell-unit-of-life:cell}{célula} y día.
\textbf{7.} $V = \frac{\pi}{6}(1)^3 = \qty{0.52}{\micro m^3} =
\qty{5.2e-13}{mL}$; proteína, $\qty{1.05e-13}{g} = \qty{0.1}{pg}$.
\textbf{8.} $125/0.1 = 1250$ gránulos al día, $0.87$ por minuto.
\textbf{9.} $V = 12^3 = \qty{1728}{\micro m^3} = \qty{1.73e-9}{cm^3}$;
masa \qty{1.8e-9}{g} = \qty{1800}{pg}. Secreta $125/1800 = 7\%$ de su masa
al día (la \omterm{def:b1:cell-unit-of-life:cell}{célula} entera es solo \qty{20}{\%} de proteína, así que un tercio
de su contenido proteico).
\textbf{10.} Masa molar $250\times 110 = \qty{27500}{g/mol}$;
$\qty{1.05e-13}{g}/27500\times\num{6e23} = \num{2.3e6}$ moléculas.
\textbf{11.} $250/5 = \qty{50}{s}$.
\textbf{12.} Moléculas al día, $1250\times\num{2.3e6} = \num{2.9e9}$; un
ribosoma fabrica $86\,400/50 = 1728$ al día; hacen falta
$\num{1.7e6}$ ribosomas, alrededor del \qty{40}{\%} de los de la \omterm{def:b1:cell-unit-of-life:cell}{célula}.
\textbf{13.} $\pi d^2 = \qty{3.14}{\micro m^2}$.
\textbf{14.} $1250\times 3.14 = \qty{3900}{\micro m^2}$.
\textbf{15.} $3900/144 = 27$ veces: la \omterm{def:b1:cell-unit-of-life:cell}{célula} recupera esa misma área por
\omterm{def:b1:eukaryotic-cell:endocytosis}{endocitosis} y la devuelve al \omterm{def:b1:eukaryotic-cell:endomembrane}{Golgi}.
\textbf{16.} $20\,000/3900 = 5$ días: el RE solo podría abastecer cinco días
de membrana de gránulos si no volviera nada; la membrana debe reciclarse sin
cesar.
\textbf{17.} Cada microvellosidad, $\pi\times 0.1\times 1 =
\qty{0.31}{\micro m^2}$; \qty{630}{\micro m^2} para 2000; factor $5.4$.
\textbf{18.} Al principio de la cadena: la \omterm{prop:b1:eukaryotic-cell:secretory}{secuencia señal}, veinte residuos
en su mayoría hidrófobos.
\textbf{19.} La señal dirigió la cadena naciente al interior de las
vesículas, donde una peptidasa señal la eliminó; la membrana protege a la
proteína de la proteasa.
\textbf{20.} Solo entran en el RE las proteínas con señal; una proteína
citosólica carece de ella y se queda fuera.
\textbf{21.} La enzima se fabrica en ribosomas libres y se queda en el
\omterm{def:b1:eukaryotic-cell:organelle}{citosol}; al ser una enzima digestiva (como precursor inactivo hace poco
daño, pero si se activara) digeriría la \omterm{def:b1:cell-unit-of-life:cell}{célula} desde dentro.
\textbf{22.} En el \omterm{def:b1:eukaryotic-cell:endomembrane}{Golgi} trans: la enzima lisosómica lleva una marca
distintiva (un azúcar fosforilado) que reconoce un receptor que la empaqueta
en vesículas destinadas a los \omterm{def:b1:eukaryotic-cell:endomembrane}{lisosomas} y no en gránulos de secreción.
\textbf{23.} En el RE (y en el \omterm{def:b1:eukaryotic-cell:endomembrane}{Golgi}): la formación de vesículas, el
transporte por motores a lo largo de los \omterm{def:b1:eukaryotic-cell:cytoskeleton}{microtúbulos} y la fusión de
membranas consumen ATP; y la propia síntesis de proteínas se detiene, de
modo que el pulso ni siquiera se completa.
\textbf{24.} \qty{45}{min} de tránsito a $0.87$ gránulos por minuto: unos
\qty{40}{} gránulos de proteína en tránsito.
\textbf{25.} Tránsito de la síntesis al gránulo, unos \qty{45}{min}
(\qty{55}{min} contando el RE); secreción al cabo de unas dos horas en el
experimento (cuatro horas o más en la glándula viva, ya que los gránulos
esperan una comida); \qty{0.9}{} gránulos liberados por minuto; el
\qty{7}{\%} de la masa de la \omterm{def:b1:cell-unit-of-life:cell}{célula} al día.
\end{solution}
